\documentclass[oneside,12pt]{Classes/aesm_edspia}

\usepackage{minitoc}
\usepackage[latin1]{inputenc}
\usepackage[english,french]{babel}
\usepackage[T1]{fontenc}
\usepackage{amsmath}
\usepackage{lmodern}%font modern
\rmfamily
\DeclareFontShape{T1}{lmr}{bx}{sc}{<->ssub * cmr/bx/sc}{} 
\usepackage{lettrine}
\usepackage{tabularx}
\usepackage{epsfig, floatflt, amssymb} 
\usepackage{moreverb} %% pour le verbatim en boite
\usepackage{cases}%equations en systemes num�rot�s - soluce possible package : CASES
\usepackage{multirow} %% pour regrouper un texte sur plusieurs lignes dans une table
\usepackage{url} %% pour citer les url par \url
\usepackage[all]{xy} %% pour la barre au dessus des symboles
\usepackage{textcomp} %% pour le symbol pour mille par \textperthousand et degr�s par \degres
\usepackage[right]{eurosym}
\usepackage{setspace} %interligne simple, double etc...
\usepackage{Classes/eurosans} %%pour le symbole \euro
\usepackage{epic,eepic}
\usepackage{soul}
\usepackage[nottoc]{tocbibind} % tables des figures, des matieres et autres dans la TOC
\usepackage{fancybox}
\usepackage[leftcaption]{sidecap}
\usepackage[labelsep=endash, textfont={footnotesize, singlespacing}, margin=10pt, format=plain, labelfont=bf]{caption}
\usepackage[Conny]{Classes/fncychap} %en tete chapitrage
\newcommand{\ie}{c.-\`a-d.~}
\hbadness=10000% pb d'overfull box r�gl�
\hfuzz=50pt
\pdfcompresslevel9 % pour compresser le pdf final au maximum
\pdfoptionpdfminorversion=5 % pour accept� les images PDF version 1.5 (ex: celles produites par Office 2007)
\def\underscore{\char`\_}
\makeatletter
\renewcommand{\thesection}{\arabic {section}}
\renewcommand{\SC@figure@vpos}{c}% centrer verticalement le caption avec le package sidecap...
\renewcommand{\fnum@figure}{\small\textbf{Figure~\thefigure}}
\renewcommand{\fnum@table}{\small\textbf{Tableau~\thetable}}

\makeatother
\usepackage{subfig}
\def\thechapter{\Roman{chapter}}

%\usepackage[framed,numbered,autolinebreaks,useliterate]{Classes/mcode}


%%% Listings

\usepackage{listings}
\lstloadlanguages{xml, java}
	
	 \usepackage{listings}
  \usepackage{courier}
 \lstset{
         basicstyle=\footnotesize\ttfamily, 
         %numbers=left,               
         numberstyle=\tiny,          
         %stepnumber=2,               
         numbersep=5pt,              
         tabsize=2,                  
         extendedchars=true,         
         breaklines=true,            
         keywordstyle=\color[rgb]{0.43,0,0}\textbf,
    		frame=b,
         commentstyle=\color[rgb]{0.51,0.51,0.51} \textit ,
         stringstyle=\ttfamily  \color[rgb]{0,0.44,0} ,
         showspaces=false,           
         showtabs=false,             
         xleftmargin=17pt,
         framexleftmargin=17pt,
         framexrightmargin=5pt,
         framexbottommargin=4pt,
         %backgroundcolor=\color{lightgray},
         showstringspaces=false            
 }
 
 \usepackage{caption}
\DeclareCaptionFont{white}{\color{white}}
\DeclareCaptionFont{red}{\color{red}}
\DeclareCaptionFont{black}{\color{black}}
\DeclareCaptionFormat{listing}{\colorbox[cmyk]{0.43, 0.35, 0.35,0.01}{\parbox{\textwidth}{\hspace{15pt}#1#2#3}}}
\captionsetup[lstlisting]{format=listing,labelfont=black,textfont=white, singlelinecheck=false, margin=0pt, font={bf,footnotesize}}


%%%%%%%%%%%%%%%%%%%%%%%%%%%%%%%%%%%%%%%%%%%
\begin{document}
%%%%%%%%%%%%%%%%%%%%%%%%%%%%%%%%%%%%%%%%%%%
\renewcommand\figurename{\small\textbf{Figure}} 

\addtocounter{page}{-1}%pour revenir � 0

% Pour remplir la page de garde
\AuteurA{Flen} {FOULENI} 
%\AuteurB{Flen2} {FOULENI}
%\AuteurC{Flen3} {FOULENI} 
%\AuteurD{Flen4} {FOULENI}

\Encadrant{Mr}{Flen}{FOULENI}
\EncadrantS{Mme} {Flena} {FALTEN}

\Filiere{GL ou RT}
\datesout{--/--/2015}



\President{M. President} {FLEN}     %% Pr�sident du Jury
\RapporteurA{Mme. Rapporteur} {FLENA} %%Rapporteur



\AnneeUniv{2014/2015}

%%%%%%%%%%%%%%%%%%%%%%%%%%%%%%%%%%%%%%%%%%%
\makethese %% cr�e la couverture.

\onehalfspacing

% une page blanche (deuxi�me de couverture)
\newpage\thispagestyle{empty}\addtocounter{page}{-3}
\null\newpage\thispagestyle{empty}


\frontmatter %num�rotation en iii
\pagestyle{fancy}
\fancyhf{}
\fancyhead[R]{Remerciements}
\fancyfoot[R]{\thepage}
\renewcommand{\headrulewidth}{0.5pt}
\renewcommand{\footrulewidth}{0pt}

\chapter*{Remerciements}
%===================================================================

Merci � tous ! Bonne journ�e


%%%%%%%% TOC

%profondeur dans la table des mati�res et de la num�rotation des sections

\setcounter{secnumdepth}{3}
\setcounter{tocdepth}{3}


\renewcommand{\contentsname}%
    {Table des Mati�res}%

%%%%minitoc
\dominitoc % g�n�re la minitoc
\nomtcrule % supprime les lignes horizontales de la minitoc
\renewcommand{\mtctitle}{Plan} % Modifie le titre de la minitoc

%%%%
\tableofcontents

\renewcommand{\headrulewidth}{0.5pt}
\renewcommand{\footrulewidth}{0pt}
\fancyhead[R]{Table des Mati�res}


%%%%%%%% Figures

\makeatletter
%\renewcommand{a\thefigure}{\@arabic\c@figure}
\@addtoreset{figure}{chapter}
\makeatother

\renewcommand{\headrulewidth}{0.5pt}
\renewcommand{\footrulewidth}{0pt}
\renewcommand\listfigurename{Liste des Figures}
\listoffigures \mtcaddchapter 

\fancyhead[R]{Liste des Figures}
\newpage


%%%%%%%% Tableaux

\makeatletter

\renewcommand{\headrulewidth}{0.5pt}
\renewcommand{\footrulewidth}{0pt}
\renewcommand\listtablename{Liste des Tableaux}

\listoftables  \mtcaddchapter 

\fancyhead[R]{Liste des Tableaux}

%%%%%%%%%%%%%%%%%%%%%%%%%%%%%%%%%%%
%\fancyhead[R]{R�sum�s}

\chapter*{R�sum�}
\addcontentsline{toc}{chapter}{R�sum�}
%===================================================================

Ceci est le r�sum� en fran�ais de votre projet. Il devra �tre plus d�taill� que le r�sum� se trouvant dans le verso de votre rapport.

\chapter*{Abstract}
\addcontentsline{toc}{chapter}{Abstract}
%===================================================================

This is the english abstract of your project. It must be longer and presented in more details than the abstract you write on the back of your report.


%%%%%%%%%%%%%%%%%%%%%%%%%%%%%%%%%%%

                       
\mainmatter %num�ros arabes
\pagestyle{fancy}
\fancyhead[R]{Introduction G�n�rale}
\chapter*{Introduction G�n�rale}

\addcontentsline{toc}{chapter}{Introduction G�n�rale}
\begin{spacing}{1.2}
%==================================================================================================%

Pour �crire un bon rapport \cite{SFAXI2015} de projet en informatique, il existe certaines r�gles � respecter. Certes, chacun �crit son rapport avec sa propre plume et sa propre signature, mais certaines r�gles restent universelles    \cite{Latex}.\\

\textbf{La Table de mati�re} est la premi�re chose qu'un rapporteur va lire. Il faut qu'elle soit :
\begin{itemize}
\item Assez d�taill�e \footnote{Sans l'�tre trop}. En g�n�ral, 3 niveaux de num�ros suffisent;
\item Votre rapport doit �tre r�parti en chapitres �quilibr�s, � part l'introduction et la conclusion, naturellement plus courts que les autres;
\item Vos titres doivent �tre suffisamment personnalis�s pour donner une id�e sur votre travail. �viter le : � Conception �,  mais privil�gier : � Conception de l'application de gestion des $...$ � M�me s'ils vous paraissent longs, c'est mieux que 
d'avoir un sommaire impersonnel. \\
\end{itemize}

\textbf{Une introduction} doit �tre r�dig�e sous forme de paragraphes bien ficel�s. Elle est
normalement constitu�e de 4 grandes parties :
\begin{enumerate}
\item Le contexte de votre application : le domaine en g�n�ral, par exemple le domaine du web, de BI, des logiciels de gestion ?
\item La probl�matique : quels sont les besoins qui, dans ce contexte l�, n�cessitent la r�alisation de votre projet?
\item La contribution : expliquer assez bri�vement en quoi consiste votre application, sans entrer dans les d�tails de r�alisation. Ne pas oublier qu'une introduction est
 cens�e introduire le travail, pas le r�sumer; 
 \item La composition du rapport : les diff�rents chapitres et leur composition. Il n'est pas n�cessaire de num�roter ces parties, mais les mettre plut�t sous forme de paragraphes successifs bien li�s.
\end{enumerate}






\end{spacing}



\fancyhf{}
\fancyhead[R]{Introduction G�n�rale}
\fancyfoot[R]{\thepage}
\renewcommand{\headrulewidth}{0.5pt}
\renewcommand{\footrulewidth}{0pt}



\setcounter{mtc}{5} %indique le num�ro r�el du chapitre, pour la mini table des mati�res
\chapter{�tude Th�orique}
\minitoc  %insert la minitoc

\graphicspath{{Chapitre1/figures/}}
%==============================================================================
\pagestyle{fancy}
\fancyhf{}
\fancyhead[R]{\bfseries\rightmark}
\fancyfoot[R]{\thepage}
\renewcommand{\headrulewidth}{0.5pt}
\renewcommand{\footrulewidth}{0pt}
\renewcommand{\chaptermark}[1]{\markboth{\MakeUppercase{\chaptername~\thechapter. #1 }}{}}
\renewcommand{\sectionmark}[1]{\markright{\thechapter.\thesection~ #1}}

\begin{spacing}{1.2}
%==============================================================================

\section*{Introduction}
Une �tude th�orique \cite{YOUSFI2015} peut contenir l'une et/ou l'autre de ces deux parties :
\section{�tat de l'art} 
C'est une �tude assez d�taill�e sur ce qui existe sur le march� ou dans la litt�rature (d'o� 
le terme �tat de l'art), qui permet de r�pondre � la probl�matique. L'id�e ici est de faire 
un comparatif entre les solutions existantes, mais surtout d'analyser le r�sultat de cette 
comparaison et de dire pourquoi ne sont-elles pas satisfaisantes pour r�pondre � votre 
probl�matique.
\begin{figure}[!ht]\centering
\includegraphics[scale=0.9]{art.jpg}
\caption{�tat de l'art}
\label{fig:fig1}
\end{figure}
\section{�tude de l'existant}
Elle est en g�n�ral r�alis�e quand on va d�velopper un module suppl�mentaire sur un 
logiciel existant, ou si on va modifier une application existante. L'�tude de l'existant
consiste � expliquer ce qui existe d�j� dans votre environnement de travail.

\section*{Conclusion}
La conclusion est en g�n�ral sans num�rotation, et n'appara�t pas dans la table des mati�res.


%==============================================================================
\end{spacing}


\setcounter{chapter}{1}
\chapter{Conception}
\minitoc %insert la minitoc
\graphicspath{{Chapitre2/figures/}}

%\DoPToC

%==============================================================================
\pagestyle{fancy}
\fancyhf{}
\fancyhead[R]{\bfseries\rightmark}
\fancyfoot[R]{\thepage}
\renewcommand{\headrulewidth}{0.5pt}
\renewcommand{\footrulewidth}{0pt}
\renewcommand{\chaptermark}[1]{\markboth{\MakeUppercase{\chaptername~\thechapter. #1 }}{}}
\renewcommand{\sectionmark}[1]{\markright{\thechapter.\thesection~ #1}}

\begin{spacing}{1.2}
%==============================================================================
\section*{Introduction}
La partie conception de l'application \textbf{\textsl{n'est pas toujours obligatoire}}. En effet, quand notre travail consiste en une �tude th�orique, ou une mise en place d'un syst�me par exemple, 
il est inutile voire obsol�te de faire un diagramme de classes ou de s�quence.\\
Quand il s'agit de d�veloppement, par contre, la partie conception s'impose. 
\section{Recommadations}
En g�n�ral, 
il faut suivre les r�gles suivantes :
\begin{itemize}
	\item Choisir une m�thodologie de travail : un processus unifi�, une m�thode agile;
\end{itemize}
\section{Diagrammes}
Il faut Bien choisir les diagrammes ad�quats pour votre application. En g�n�ral, les
diagrammes obligatoires sont les diagrammes de cas d'utilisation, de classe et de 
s�quence. Vous pouvez ajouter en plus le diagramme qui vous semble pertinent :
par exemple, pour une application sur plusieurs tiers, il est int�ressant de
montrer le diagramme de d�ploiement;
\begin{itemize}

\item Les diagrammes doivent �tre clairs, lisibles et bien expliqu�s, sans pour autant
nous submerger de d�tails. Des explications trop longues deviennent ennuyeuses;
\item Si un diagramme est trop grand, vous pouvez le diviser, le repr�senter sous
forme de plusieurs diagrammes, ou vous abstraire de certains d�tails. Si c'est
impossible, imprimez le sur une grande page (A3), quitte � la plier ensuite. Le
plus important est que tous les mots soient lisibles.

\end{itemize}
\subsection{Diagramme de S�quence}
Un diagramme de s�quence :
\begin{itemize}
	\item Repr�sente un sc�nario possible qui se d�roule dans un cas d'utilisation. 
	Vous n'�tes donc pas oblig�s de montrer tous les cas d'ex�cution possibles;
	\item Repr�sente l'interaction entre les objets : donc normalement, toutes les
	instances d�finies dans un diagramme de s�quences doivent correspondre 
	 � des classes qui se trouvent dans le diagramme des classes;
	 \item Il existe parfois des dizaines de diagrammes de s�quences possibles. Choisissez certains d'entre eux � mettre dans le rapport (2 ou 3). Privil�giez les diagrammes les plus importants (et non, l'authentification n'en fait pas partie!).
\end{itemize}
\subsection{Diagramme de Classes}
Un diagramme de classes :
\begin{itemize}
\item Doit �tre fid�le � l'architecture logicielle choisie. Si vous utilisez le MVC, 
alors les trois couches doivent �tre repr�sent�es dans le diagramme de classes gr�ce aux packages;
\item Les st�r�otypes sont fortement conseill�s. Si vous d�veloppez une
application web, n'h�sitez pas � utiliser les st�r�otypes de la figure \ref{fig:fig2} : 
\begin{figure}[!ht]\centering
\includegraphics[scale=0.9]{stereotypes.jpg}
\caption{Les st�r�otypes}
\label{fig:fig2}
\end{figure}
\item Attention � ne pas confondre classes et tables : �vitez la tentation de
mettre des id partout.

\end{itemize}
 
\section*{Conclusion}
Faire ici une petite r�capitulation du chapitre, ainsi qu'une introduction du chapitre suivant.





%==============================================================================
\end{spacing}


\setcounter{chapter}{2}
\chapter{R�alisation}
\minitoc %insert la minitoc
\graphicspath{{Chapitre3/figures/}}

%\DoPToC
%==============================================================================
\pagestyle{fancy}
\fancyhf{}
\fancyhead[R]{\bfseries\rightmark}
\fancyfoot[R]{\thepage}
\renewcommand{\headrulewidth}{0.5pt}
\renewcommand{\footrulewidth}{0pt}
\renewcommand{\chaptermark}[1]{\markboth{\MakeUppercase{\chaptername~\thechapter. #1 }}{}}
\renewcommand{\sectionmark}[1]{\markright{\thechapter.\thesection~ #1}}

\begin{spacing}{1.2}

%==============================================================================
\section*{Introduction}
Ce chapitre porte sur la partie pratique ainsi que la bibliographie.

\section{Outils et langages utilis�s}
L'�tude technique peut se trouver dans cette partie, comme elle peut �tre faite en
parall�le avec l'�tude th�orique (comme le sugg�re le mod�le 2TUP).
Dans cette partie, il faut essayer de convaincre le lecteur de vos choix en termes de
technologie. Un �tat de l'art est souhait� ici, avec un comparatif, une synth�se et un choix 
d'outils, m�me tr�s brefs.
\section{Pr�sentation de l'application}
Il est tout � fait normal que tout le monde attende cette partie pour coller � souhait toutes les images
correspondant aux interfaces diverses de l'application si ch�re � votre coeur, mais
abstenez vous! Il FAUT mettre des imprime �crans, mais bien choisis, et surtout, il faut les sc�nariser : Choisissez un sc�nario d'ex�cution, par exemple la cr�ation d'un 
nouveau client, et montrer les diff�rentes interfaces n�cessaires pour le faire, en
expliquant bri�vement le comportement de l'application. Pas trop d'images, ni trop de
commentaires : concis, encore et toujours.

�vitez ici de coller du code : personne n'a envie de voir le contenu de vos classes.
Mais  vous  pouvez ins�rer des snippets (bouts de code) pour montrer certaines
fonctionnalit�s \cite{YOUSFI2015}\cite{Latex}, si vous en avez vraiment besoin. Si vous voulez montrer une partie de votre code, les �tapes d'installation ou de configuration, vous pourrez les mettre dans l'annexe.
\subsection{Exemple de tableau}

Vous pouvez utiliser une description tabulaire d'une �ventuelle comparaison entre les travaux existants. Ceci est un exemple de tableau: Tab \ref{tab:exple}.

\begin{table}[ht]
	\centering
	\caption{Tableau comparatif}
	\footnotesize
	\begin{tabularx}{\linewidth}{|>{\bfseries \vspace*{\fill}}X ||>{\centering{}\vspace*{\fill}}X|>{\centering{}\vspace*{\fill}}X|>{\centering{}\vspace*{\fill}}X|>{\vspace*{\fill}}X<{\centering{}}|}	
			\hline 
			& \bfseries Col1 & \bfseries Col2 &\bfseries Col3 &\bfseries Col4\\
			\hline \hline
			Row1		&		&	X	&		&		\\
			Row2		&	X	&		&		&		\\
			Row3		&	X	&	X	&	X	&	X	\\
			Row4		&	X	&		&	X	&	X	\\
			Row5		&	X	&		&	X	&	X	\\
			Row6		&	X	&		&	X	&	X	\\
			Row7		&	X	&		&	X	&		\\
			Row8		&	X	&	X	&	X	&		\\
			\hline
	\end{tabularx}
	\label{tab:exple}
\end{table}

\subsection{Exemple de Code}
Voici un exemple de code Java, avec coloration syntaxique \ref{code:java}.

\begin{lstlisting}[rulecolor=\color{white}]
\end{lstlisting}

\begin{lstlisting}[label=code:java,caption=Helloworld Java,language=java]
	public class HelloWorld {
	//la m�thode main
    public static void main(String[] args) {
        System.out.println("Hello, World");
    }

}
\end{lstlisting}

\section{Remarques sur la bibliographie}
Votre bibliographie doit r�pondre � certains crit�res, sinon, on vous fera encore et
toujours la remarque dessus (et parfois, m�me si vous pensez avoir tout fait comme il
 faut, on peut vous faire la remarque quand m�me : chacun a une conception tr�s
personnelle de comment une bibliographie devrait �tre).\\
\begin{itemize}
\item Une bibliographie dans un bon rapport doit contenir plus de livres et d'articles 
que de sites web : apr�s tout c'est une biblio. Privil�giez donc les ouvrages
reconnus et publi�s pour vos d�finitions, au lieu de sauter directement sur le premier article wikipedia;
 \item Les �l�ments d'une bibliographie sont de pr�f�rence class�s par ordre
alphab�tique, ou par th�mes (et ordre alphab�tique pour chaque th�me);
\item Une entr�e bibliographique doit �tre sous la forme suivante :
\begin{itemize}
\item Elle doit contenir un identifiant unique: repr�sent� soit par un num�ro
[1] ou par le nom du premier auteur, suivi de l'ann�e d'�dition [Kuntz, 1987];
\item Si c'est un livre : Les noms des auteurs, suivi du titre du livre, de l'�diteur, 
ISBN/ISSN, et la date d'�dition;
\item Si c'est un article : Les noms des auteurs, le titre , le journal ou la
conf�rence, et la date de publication;
\item Si c'est un site web ou un document �lectronique : Le titre, le lien et la date 
de consultation;
\item Si c'est une th�se : nom et pr�nom, titre de la th�se, universit� de
soutenance, ann�e de soutenance, nombre de pages;
\item Exemples : 
\begin{description}
\item $[Bazin, 1992]$ BAZIN R., REGNIER B. Les traitements antiviraux et leurs essais
th�rapeutiques. Rev. Prat., 1992, 42, 2, p.148-153.\\
\item $[Anderson,1998]$ ANDERSON P.JF. Checklist of criteria used for evaluation of metasites.
[en ligne]. Universit� du Michigan, Etats Unis. Site disponible sur :\\
http://www.lib.umich.edu/megasite/critlist.html.(Page consult�e le 11/09/1998).
\end{description}
\item Dans le texte du rapport, on doit obligatoirement citer la r�f�rence en  faisant appel � son identifiant, juste apr�s avoir utilis� la citation. Si ceci n'est pas fait dans les r�gles, on peut �tre accus� de plagiat.
\end{itemize} 
\end{itemize} 

\section*{Conclusion}
Voil�.

%==============================================================================
\end{spacing}


\backmatter
\pagestyle{fancy}
\fancyhf{}
\renewcommand{\chaptermark}[1]{\markboth{Conclusion G�n�rale et Perspectives}{}}
\fancyhead[R]{Conclusion G�n�rale et Perspectives}
\fancyfoot[R]{\thepage}
\renewcommand{\headrulewidth}{0.5pt}
\renewcommand{\footrulewidth}{0pt}
\chapter{Conclusion G�n�rale et Perspectives}
%==============================================================================
\pagestyle{fancy}
\fancyhf{}
\fancyhead[R]{\bfseries\rightmark}
\fancyfoot[R]{\thepage}
\renewcommand{\headrulewidth}{0.5pt}
\renewcommand{\footrulewidth}{0pt}
\renewcommand{\chaptermark}[1]{\markboth{\MakeUppercase{\chaptername~\thechapter. #1 }}{}}
\renewcommand{\sectionmark}[1]{\markright{\thechapter.\thesection~ #1}}

\begin{spacing}{1.2}
%==============================================================================

C'est l'une des parties les plus importantes et pourtant les plus n�glig�es 
du rapport. Ce qu'on \underline{ne veut pas voir} ici, c'est combien ce stage vous a �t� b�n�fique, comment il vous a appris � vous int�grer, � conna�tre le monde du travail, etc.\\
Franchement, personne n'en a rien � faire, du moins dans cette partie. Pour cela, vous 
avez les remerciements et les d�dicaces, vous pourrez vous y exprimer � souhait.\\
La conclusion, c'est tr�s simple : c'est d'abord le r�sum� de ce que vous avez racont�
dans le rapport : vous reprenez votre contribution, en y ajoutant ici les outils que vous 
avez utilis�, votre mani�re de proc�der. Vous pouvez m�me mettre les difficult�s
rencontr�es. En deuxi�me lieu, on y met les perspectives du travail : ce qu'on pourrait 
ajouter � votre application, comment on pourrait l'am�liorer.

%==============================================================================
\end{spacing}

\bibliographystyle{Biblio/unsrt_modif} 
\singlespacing
\renewcommand{\bibname}{Bibliographique}

\bibliography{Biblio/aesm_edspia}

\onehalfspacing

\appendix
\setcounter{figure}{0} 
\setcounter{table}{0}
\setcounter{footnote}{0}
\setcounter{equation}{0}
\pagestyle{fancy}
\fancyhf{}
\renewcommand{\chaptermark}[1]{\markboth{\MakeUppercase{#1 }}{}}
\renewcommand{\sectionmark}[1]{\markright{\thesection~ #1}}
\fancyhead[RO]{\bfseries\rightmark}
\fancyhead[LE]{\bfseries\leftmark}
\fancyfoot[RO]{\thepage}
\fancyfoot[LE]{\thepage}
\renewcommand{\headrulewidth}{0.5pt}
\renewcommand{\footrulewidth}{0pt}

\makeatletter
\renewcommand\thefigure{A.\arabic{figure}}
\renewcommand\thetable{A.\arabic{table}} 
\makeatother

\chapter{Annexe : Remarques Diverses}
\graphicspath{{Annexe1/figures/}}
%==========================================================================

%    Annexe

%===========================================================================
\begin{itemize}
\item Un rapport doit toujours �tre bien num�rot�;
\item De pr�f�rence, ne pas utiliser plus que deux couleurs, ni un caract�re fantaisiste; 
\item Essayer de toujours garder votre rapport sobre et professionnel; 
\item Ne jamais utiliser de je ni de on, mais toujours le nous (m�me si tu as tout fait tout seul); 
\item Si on n'a pas de paragraphe 1.2, ne pas mettre de 1.1;
\item TOUJOURS, TOUJOURS faire relire votre rapport � quelqu'un d'autre (de pr�f�rence qui n'est pas du domaine) pour vous corriger les fautes d'orthographe et de fran�ais;
\item Toujours valoriser votre travail : votre contribution doit �tre bien claire et mise en �vidence; 
\item Dans chaque chapitre, on doit trouver une introduction et une conclusion;
\item Ayez toujours un fil conducteur dans votre rapport. Il faut que le lecteur suive un raisonnement bien clair, et trouve la relation entre les diff�rentes parties;
\item Il faut toujours que les abr�viations soient d�finies au moins la premi�re fois o� elles sont utilis�es. Si vous en avez beaucoup, utilisez un glossaire.
\item Vous avez tendance, en d�crivant  l'environnement mat�riel, � parler de votre ordinateur, sur lequel vous avez d�velopp� : ceci est inutile. Dans cette partie, on ne cite que le mat�riel qui a une influence sur votre application. Que vous l'ayez d�velopp� sur Windows Vista ou sur Ubuntu n'a aucune importance;
\item Ne jamais mettre de titres en fin de page; 
\item Essayer toujours d'utiliser des termes fran�ais, et �viter l'anglicisme. Si certains termes  sont plus connus en  anglais, donner leur �quivalent en fran�ais la premi�re fois que vous les utilisez, puis utilisez le mot anglais, mais en italique;
\item �viter les phrases trop longues : clair et concis, c'est la r�gle g�n�rale !\\

\newpage

\textbf{Rappelez vous que votre rapport est le visage de votre travail : un mauvais rapport peut �clipser de l'excellent travail. Alors pr�tez-y l'attention n�cessaire.}

 
\begin{figure}[!ht]\centering
\includegraphics[scale=0.5]{ingenieur.jpg}
\end{figure}
\end{itemize}



\end{document}